\subsection{环的变换}

命题 5.—— 设 \(\rho: A \to B\) 为 Noether 环的同态，它使 B 成为平坦 A-模。假设对 B 的每个极大理想 \(\mathfrak{n}\)，环 \(\kappa(\rho^{-1}(\mathfrak{n})) \otimes_{A} B\) 都是 Gorenstein 环。设 \(\Omega\) 为对偶化 A-模；则 B-模 \(\Omega_{(B)}\) 是对偶化的。

设 \(\mathfrak{n}\) 为 B 的一个极大理想，并用 \(\mathfrak{p}\) 表示它在 A 中的逆像。\(A_{\mathfrak{p}}\)-模 \(B_{\mathfrak{n}}\) 是平坦的，\(A_{\mathfrak{p}}\)-模 \(\Omega_{\mathfrak{p}}\) 是对偶化的，\(\Omega_{(B)} \otimes_B B_{\mathfrak{n}}\) 等同于 \(\Omega_{\mathfrak{p}} \otimes_{A_{\mathfrak{p}}} B_{\mathfrak{n}}\)，而 \(\kappa_{\Lambda_{\mathfrak{p}}} \otimes_{\Lambda_{\mathfrak{p}}} B_{\mathfrak{n}}\)——它等同于 \(\kappa(\mathfrak{p}) \otimes_A B\) 的一个分式环——是 Gorenstein 环。因此，只需在 \(\rho\) 是局部环之间局部同态的情形证明命题，以下即作此假设。

先处理环 A 与 B 都是 Artin 环的情形。令 \(C = B / \mathfrak{m}_{\Lambda}B\)。由于 B 在 A 上平坦，B-模 \(\mathrm{Hom}_{\mathrm{B}}(\mathrm{C},\Omega_{(\mathrm{B})})\) 同构于 \(\mathrm{Hom}_{\Lambda}(\kappa_{\mathrm{A}},\Omega)\otimes_{\mathrm{A}}\mathrm{B}\)（I, § 2, 第 10 节，命题 11），从而同构于 \(\mathrm{Hom}_{\mathrm{A}}(\kappa_{\mathrm{A}},\Omega)\otimes_{\kappa_{\mathrm{A}}}\mathrm{C}\)。由此推出一列同构

\[
\begin{array}{c} \operatorname{Hom} _ {\mathrm{B}} (\kappa_ {\mathrm{B}}, \Omega_ {(\mathrm{B})}) \longrightarrow \operatorname{Hom} _ {\mathrm{C}} (\kappa_ {\mathrm{C}}, \operatorname{Hom} _ {\mathrm{B}} (\mathrm{C}, \Omega_ {(\mathrm{B})})) \longrightarrow \operatorname{Hom} _ {\mathrm{C}} (\kappa_ {\mathrm{C}}, \operatorname{Hom} _ {\mathrm{A}} (\kappa_ {\mathrm{A}}, \Omega) \otimes_ {\kappa_ {\mathrm{A}}} \mathrm{C}) \\ \qquad \qquad \qquad \qquad \qquad \qquad \qquad \qquad \qquad \qquad \longrightarrow \operatorname{Hom} _ {\mathrm{A}} (\kappa_ {\mathrm{A}}, \Omega) \otimes_ {\kappa_ {\mathrm{A}}} \operatorname{Hom} _ {\mathrm{C}} (\kappa_ {\mathrm{C}}, \mathrm{C}). \end{array}
\]

\(\kappa_{\mathrm{A}}\)-向量空间 \(\mathrm{Hom}_{\mathrm{A}}(\kappa_{\mathrm{A}}, \Omega)\) 的维数为 1，因为 \(\Omega\) 是对偶化的；又由于 C 是 Gorenstein 环，\(\kappa_{\mathrm{C}}\)-向量空间 \(\mathrm{Hom}_{\mathrm{C}}(\kappa_{\mathrm{C}}, \mathrm{C})\) 同样如此；于是 \(\kappa_{\mathrm{B}}\)-向量空间 \(\mathrm{Hom}_{\mathrm{B}}(\kappa_{\mathrm{B}}, \Omega_{(\mathrm{B})})\) 的维数为 1。

设 M 为有限长度的 B-模；我们对 \(\operatorname{long}_{\mathrm{B}}(\mathrm{M})\) 作归纳，证明 \(\operatorname{long}_{\mathrm{B}}(\operatorname{Hom}_{\mathrm{B}}(\mathrm{M},\Omega_{(\mathrm{B})})) \leqslant \operatorname{long}_{\mathrm{B}}(\mathrm{M})\)。当 M = 0 时断言显然，而当 \(M = \kappa_{B}\) 时由前所述可得。设 \(\operatorname{long}_{\mathrm{B}}(\mathrm{M}) \geqslant 2\)。存在 B-模的一个正合列

\[
0 \rightarrow \mathrm{M} ^ {\prime} \rightarrow \mathrm{M} \rightarrow \kappa_ {\mathrm{B}} \rightarrow 0
\]

其中 \(\mathrm{long}_{\mathrm{B}}(\mathrm{M}^{\prime}) < \mathrm{long}_{\mathrm{B}}(\mathrm{M})\)。由此导出一个正合列

\[
0 \to \operatorname{Hom} _ {\mathrm{B}} (\kappa_ {\mathrm{B}}, \Omega_ {(\mathrm{B})}) \to \operatorname{Hom} _ {\mathrm{B}} (\mathrm{M}, \Omega_ {(\mathrm{B})}) \to \operatorname{Hom} _ {\mathrm{B}} (\mathrm{M} ^ {\prime}, \Omega_ {(\mathrm{B})}),
\]

然后对 M′ 应用归纳假设即得结论。

设 N 为从 \(\kappa_{A} \otimes_{A} B\) 到 \(\kappa_{B}\) 的典范满射的核。令 \(m = \text{long}_{\text{B}}(\kappa_{\text{A}} \otimes_{\text{A}} \text{B})\)；则有 \(\text{long}_{\text{B}}(\text{N}) = m - 1\)。考虑 B-模的正合列

\[
\begin{array}{c} 0 \to \operatorname{Hom} _ {\mathrm{B}} (\kappa_ {\mathrm{B}}, \Omega_ {(\mathrm{B})}) \longrightarrow \operatorname{Hom} _ {\mathrm{B}} (\kappa_ {\mathrm{A}} \otimes_ {\mathrm{A}} \mathrm{B}, \Omega_ {(\mathrm{B})}) \longrightarrow \operatorname{Hom} _ {\mathrm{B}} (\mathrm{N}, \Omega_ {(\mathrm{B})}) \\ \longrightarrow \operatorname{Ext} _ {\mathrm{B}} ^ {1} (\kappa_ {\mathrm{B}}, \Omega_ {(\mathrm{B})}) \longrightarrow \operatorname{Ext} _ {\mathrm{B}} ^ {1} (\kappa_ {\mathrm{A}} \otimes_ {\mathrm{A}} \mathrm{B}, \Omega_ {(\mathrm{B})}). \end{array}
\]

B-模 \(\mathrm{Hom}_{\mathbb{B}}(\kappa_{\Lambda}\otimes_{\Lambda}\mathrm{B},\Omega_{(\mathbb{B})})\) 与 \(\mathrm{Ext}_{\mathbb{B}}^{1}(\kappa_{\Lambda}\otimes_{\Lambda}\mathrm{B},\Omega_{(\mathbb{B})})\) 分别同构于 \(\mathrm{Hom}_{\Lambda}(\kappa_{\Lambda},\Omega)\otimes_{\Lambda}\mathrm{B}\) 与 \(\mathrm{Ext}_{\mathbb{A}}^{1}(\kappa_{\Lambda},\Omega)\otimes_{\Lambda}\mathrm{B}\)，亦即分别同构于 \(\kappa_{\Lambda}\otimes_{\Lambda}\mathrm{B}\) 与 0。B-模 \(\mathrm{Hom}_{\mathbb{B}}(\kappa_{\mathbb{B}},\Omega_{(\mathbb{B})})\) 与 \(\mathrm{Hom}_{\mathbb{B}}(\kappa_{\mathbb{A}}\otimes_{\mathbb{A}}\mathrm{B},\Omega_{(\mathbb{B})})\) 的长度分别为 1 与 \(m\)，而 \(\mathrm{Hom}_{\mathbb{B}}(\mathrm{N},\Omega_{(\mathbb{B})}))\) 的长度 \(\leqslant m - 1\)；由此推出 B-模 \(\mathrm{Ext}_{\mathbb{B}}^{1}(\kappa_{\mathbb{B}},\Omega_{(\mathbb{B})})\) 为零。由 § 3, 第 3 节，命题 6，B-模 \(\Omega_{(\mathbb{B})}\) 是内射的；从而它是对偶化模（第 1 节，例 1）。

转到一般情形，令 \(C = \kappa_A \otimes_A B\)；按假设它是 Gorenstein 环，因而是 Macaulay 环（\(\S\)。由第 1 节的命题 1，A 是 Macaulay 环，且 A-模 \(\Omega\) 是 Macaulay 模。于是 B 是 Macaulay 环，且 B-模 \(\Omega_{(B)}\) 是 Macaulay 模（\(\S\) 2, 第 7 节，命题 9 的推论 1）。令 \(r = \dim(A)\)，\(s = \dim(C)\)。存在 \(x_1, \ldots, x_r\) 的元素的一个序列（\(\mathfrak{m}_A\)），它对 A-模 A 与 \(\Omega\) 都是正则的；又存在 \(y_1, \ldots, y_s\) 的元素的一个序列（\(\mathfrak{m}_B\)），它对 B-模 C 是正则的；用 \(x\) 表示上述第一个序列在 A 中生成的理想，用 \(y\) 表示上述第二个序列在 B 中生成的理想。序列（\(y_1, \ldots, y_s, \rho(x_1), \ldots, \rho(x_r)\)）对 B-模 B 与 \(\Omega_{(B)}\) 都是正则的（\(\S\) 1, 第 6 节，命题 11），且 A-模 B/\(y\) 是平坦的（前引文，命题 10）。令 \(A' = A/x\)，\(B' = B/(xB + y)\)，并用 \(\rho': A' \to B'\) 表示由 \(\rho\) 通过取商导出的同态。环 \(A'\) 与 \(B'\) 都是 Artin 环，A′-模 \(B'\) 是平坦的，环 \(\kappa_{A'} \otimes_{A'} B'\)——它等同于 C/\(y\)——是 Gorenstein 环（\(\S\) 3, 第 7 节，例 2），且 A′-模 \(\Omega_{(A')}\) 是对偶化的（第 2 节，命题 4）。由证明的第一部分，B′-模 \(\Omega_{(B')}\) 是对偶化的。于是由前引文，B-模 \(\Omega_{(B)}\) 是对偶化的。

推论.—— 设 A 为具有对偶化模 \(\Omega\) 的 Noether 环，B 为 A 上有限多个未定元的多项式代数。则 B-模 \(\Omega_{(B)}\) 是对偶化的。

事实上，对 A 的每个素理想 p，环 \(\kappa(\mathfrak{p})\otimes_{\mathrm{A}}\mathrm{A}[\mathbf{X}]\) 等同于 \(\kappa(\mathfrak{p})[\mathbf{X}]\)，后者是正则的，因而是 Gorenstein 的。

命题 6.—— 设 \(\Lambda\) 为 Noether 局部环，\(\Omega\) 为对偶化 A-模。设 B 为有限 A-代数；假设 A-模 B 是 Macaulay 模。B-模 \(\mathrm{Ext}_{\mathrm{A}}^{i}(\mathrm{B},\Omega)\) 当 \(i\neq \dim (\Lambda) - \dim (\mathrm{B})\) 时为零，而当 \(i = \dim (\Lambda) - \dim (\mathrm{B})\) 时是对偶化的。

我们有 \(\dim(B) = \dim_{A}(B) \leqslant \dim(\Lambda)\)（VIII, § 2, 第 3 节，定理 1 c)）；令 \(c = \dim(A) - \dim(B)\)。有 \(\operatorname{Ext}_{A}^{i}(B, \Omega) = 0\)（其中 \(i \neq c\)），因为 A-模 B 是 Macaulay 模（第 1 节，命题 3 的推论）。我们来证明 B-模 \(\operatorname{Ext}_{\Lambda}^{c}(B, \Omega)\) 是对偶化的。

先设 \(\dim(B) = 0\)。B 的谱 X 是有限的，且由极大理想组成（IV, § 2, 第 5 节，命题 9）；典范映射 \(B \to \prod_{n \in X} B_n\) 是同构（前引文，推论 1）。于是 B-模 \(\Omega' = \text{Ext}_A^c(B, \Omega)\) 是诸模 \(\text{Ext}_A^c(B_n, \Omega)\) 的直和；由于 \(\text{Ext}_A^c(B_n, \Omega)\) 的支撑集含于 \(\{\mathfrak{n}\}\)，它等同于 \(\Omega'_n\)。有 \(\dim(B_n) = 0\)（对每个 \(\mathfrak{n}\)）；因此，为证明 B-模 \(\Omega'\) 是对偶化的，只需证明 \(B_n\)-模 \(\text{Ext}_A^c(B_n, \Omega)\)（对每个 \(\mathfrak{n} \in X\)）都如此，这就化归为环 B 是局部环的情形。在此情形，由 § 8, 第 5 节，例 6，B-模 \(\text{Ext}_A^c(B, \Omega)\) 同构于 \(\text{Hom}_A(B, I)\)，其中 I 是 Matlis A-模；从而它是 Matlis B-模（§ 8, 第 6 节，命题 5 的推论），因此是对偶化 B-模（第 1 节，例 1）。

现设 \(\dim(B) > 0\)，并对 \(\dim(B)\) 作归纳推理。有 \(\operatorname{prof}_{A}(B) = \dim_{A}(B) = \dim(B)\)，从而 \(\operatorname{prof}_{A}(B) > 0\)；另一方面，有 \(\operatorname{prof}(A) = \dim(A) > 0\)（第 1 节，命题 1），于是 \(\operatorname{prof}_{A}(A \oplus B) > 0\)。因此存在 \(x\) 的一个元素 \(\mathfrak{m}_{\Lambda}\)，使同态 \(x_{\Lambda}\) 与 \(x_{B}\) 都是单射的。

考虑与正合列 \(0 \to \mathrm{B} \xrightarrow{x_{\mathrm{B}}} \mathrm{B} \longrightarrow \mathrm{B}/x\mathrm{B} \to 0\) 以及与 A-模 \(\Omega\) 相伴的扩张模正合列。A-模 B/xB 是 Macaulay 模（\(\S 2, n^{\circ} 1\)，例 3），维数为 \(\dim(\mathrm{B}) - 1\)（VIII, \(\S 3, n^{\circ} 2\)）；因此有 \(\operatorname{Ext}_{\mathrm{A}}^{i}(\mathrm{B}/x\mathrm{B}, \Omega) = 0\)（对 \(i \neq c + 1\)）（\(n^{\circ} 1\)，命题 3 的推论）。由于有 \(\operatorname{Ext}_{\mathrm{A}}^{i}(\mathrm{B}, \Omega) = 0\)（对 \(i \neq c\)），故得到 B-模的一个正合列

\[
0 \rightarrow \operatorname{Ext} _ {\mathrm{A}} ^ {c} (\mathrm{B}, \Omega) \xrightarrow {x} \operatorname{Ext} _ {\mathrm{A}} ^ {c} (\mathrm{B}, \Omega) \longrightarrow \operatorname{Ext} _ {\mathrm{A}} ^ {c + 1} (\mathrm{B} / x \mathrm{B}, \Omega) \rightarrow 0.
\]

由归纳假设，B/xB-模 \(\operatorname{Ext}_{\Lambda}^{c+1}(\mathrm{B}/x\mathrm{B},\Omega)\) 是对偶化的。由于 A-代数 B 是有限的，\(m_{A}\) 在 B 中的像含于 B 的根（V, § 2, 第 1 节，命题 1）；由第 2 节的命题 4，B-模 \(\operatorname{Ext}_{\mathrm{A}}^{c}(\mathrm{B},\Omega)\) 是对偶化的。

推论 1.—— 设 A 为 Noether 环，Ω 为对偶化 A-模，B 为有限 A-代数；假设 A-模 B 是 Macaulay 模。则 B-模 Ext \(_{A}\) (B, Ω) 是对偶化的。

设 \(\Omega'\) 为 B-模 \(\mathrm{Ext}_{\mathsf{A}}(\mathsf{B},\Omega)\)。设 \(\mathfrak{n}\) 为 B 的一个极大理想；它在 A 中的逆像是极大理想 \(\mathfrak{m}\)（V, § 2, 第 1 节，命题 1）。\(\Lambda_{\mathfrak{m}}\)-代数 \(\mathrm{B}_{\mathfrak{m}} = \mathrm{A}_{\mathfrak{m}} \otimes_{\mathsf{A}} \mathrm{B}\) 是有限的，且是 Macaulay \(\mathrm{A}_{\mathfrak{m}}\)-模；由命题，\(\mathrm{B}_{\mathfrak{m}}\)-模 \(\Omega_{\mathfrak{m}}'\)——它等同于 \(\mathrm{Ext}_{\mathrm{A}_{\mathfrak{m}}}(B_{\mathfrak{m}},\Omega_{\mathfrak{m}})\)（§ 3, 第 2 节，命题 2）——是对偶化的。由于 \(\mathrm{B}_{\mathfrak{n}}\) 是 \(\mathrm{B}_{\mathfrak{m}}\) 的一个分式环，\(\mathrm{B}_{\mathfrak{n}}\)-模 \(\Omega_{\mathfrak{n}}'\) 是对偶化的，推论得证。

注.—— 保留推论 1 的假设，并进一步设典范同态 \(\rho : A \to B\) 是单射的。于是有 \(\dim(\Lambda_{\mathfrak{m}}) = \dim(B_{\mathfrak{m}})\)，其中 \(\mathfrak{m}\) 取遍 \(A\) 的极大理想。由命题 6 与推论 1，\(\operatorname{Ext}_{A}^{i}(B, \Omega)\) 当 \(i \neq 0\) 时为零，且 B-模 \(\operatorname{Hom}_{A}(B, \Omega)\) 是对偶化的。

推论 2.—— 若 Noether 环 A 具有对偶化模，则每个是 Macaulay 环的有限生成 A-代数都具有对偶化模。

这由推论 1 及命题 5 的推论得出。

推论 3.—— 每个可表现的 Macaulay 环（特别地，每个完备局部 Macaulay 环）都具有对偶化模。

事实上，设 R 为正则环，A 为 R 的一个 Macaulay 商环。R-模 A 是 Macaulay 模（§ 2, 第 5 节，例 5），且 R 具有对偶化模（第 1 节，例 2）；于是由推论 1，A 亦然。此外，已经指出完备局部 Noether 环是可表现的（§ 4, 第 4 节，命题 6, c)）。

更一般地，每个是 Gorenstein 环之商的 Macaulay 环都具有对偶化模。反之，可以证明具有对偶化模的局部 Macaulay 环是局部 Gorenstein 环的商（习题 1）。

